\documentclass{article}
\usepackage{../math_template}

\author{Jonathan Kasongo}
\title{Japan Math Olympiad 1991 --- P3/5}

\begin{document}
\maketitle

\begin{problem}{Japan Math Olympiad 1991 --- P3/5}
Let \(A\) be a positive 16-digit integer. Prove that you can select
some sequence of consecutive digits from \(A\) such that their product
forms a perfect square. (A single digit that is a perfect square is
also valid).
\end{problem}

\begin{solution}{Solution}
Let $A = \overline{d_1 d_2 \cdots d_{16}}$. Since
$0 \leq d_i \leq 9$ the product of some set of digits can only have
prime factors $2,3,5$ or 7. For prime $p$, and integer $1 \leq n \leq 16$,
define
$$
f_p(n) = \nu_p (d_1 \cdot d_2 \cdots d_n) \text{ mod } 2.
$$
Note that if there exists $n$ such that
$f_2(n) = f_3(n) = f_5(n) = f_7(n) = 0$ then we are done because
$d_1 \cdot d_2 \cdots d_n$ is square since all the exponents are even in
it's prime factorisation. Suppose that such $n$ doesn't exist. Consider
the 4-tuple
$$
T_k := (f_2(k), f_3(k), f_5(k), f_7(k)), \quad (k=1,2,\ldots,16).
$$
Each $f_p(k)$ is either 0,1, and so there are $16 - 1 = 15$ possible options
for the tuple $T_k$, since we are excluding $(0,0,0,0)$. Since we have
$16 > 15$ tuples $T_1, T_2, \ldots, T_{16}$, by the Pigeonhole principle
there must be two tuples $T_i$ and $T_j$, $(i < j)$ such that
$T_i = T_j$. Now because
$$
\nu_p(d_i \cdot d_{i+1} \cdots d_j) = \nu_p(d_1 \cdot d_2 \cdots d_j) -
\nu_p(d_1 \cdot d_2 \cdots d_{i-1})
$$
and the fact that for $p=2,3,5,7$ the two quantities on the RHS are the same
parity, we deduce that $d_i \cdot d_{i+1}\cdots d_j$ is square, since it
only has even exponents in it's prime factorisation. $\blacksquare$
\end{solution}

\end{document}
